% ============================================================================
% Preámbulo común del Informe Técnico: Agente Solucionador de Laberintos 2D
% ============================================================================
\usepackage[utf8]{inputenc}
\usepackage[T1]{fontenc}
\usepackage{lmodern}
\usepackage{microtype}
\usepackage[spanish,es-tabla,es-nodecimaldot]{babel}

% Geometría y encabezados
\usepackage[letterpaper,top=2.4cm,bottom=2.4cm,left=2.4cm,right=2.4cm,headheight=15pt]{geometry}
\usepackage{fancyhdr}
\usepackage{titlesec}
\usepackage{parskip}

% Matemáticas y teoremas
\usepackage{amsmath,amssymb,amsthm}
\theoremstyle{definition}
\newtheorem{definicion}{Definición}[section]
\theoremstyle{plain}
\newtheorem{proposicion}{Proposición}[section]
\newtheorem{teorema}{Teorema}[section]

% Unidades SI
\usepackage{siunitx}
\sisetup{
  output-decimal-marker = {.},
  per-mode = symbol,
  detect-all = true,
  range-phrase = {--},
  range-units = single
}

% Colores institucionales y técnicos
\usepackage[table,dvipsnames]{xcolor}
\definecolor{unalverde}{HTML}{1B5E20}
\definecolor{azuldoc}{HTML}{154360}
\definecolor{azulclaro}{HTML}{2E86C1}
\definecolor{grisborde}{HTML}{BDC3C7}
\definecolor{gristexto}{HTML}{4D5656}
\definecolor{ambarborde}{HTML}{B7791F}
\definecolor{ambarfondo}{HTML}{FEF9E7}
\definecolor{azulfondo}{HTML}{EBF5FB}
\definecolor{verdefondo}{HTML}{EAFAF1}
\definecolor{rojoborde}{HTML}{A93226}
\definecolor{rojofondo}{HTML}{FDEDEC}
\definecolor{codigofondo}{HTML}{F8F9F9}

% Figuras, tablas y listas
\usepackage{graphicx}
\usepackage{float}
\usepackage{caption}
\usepackage{subcaption}
\usepackage{booktabs}
\usepackage{tabularx}
\usepackage{longtable}
\usepackage{multirow}
\usepackage{array}
\usepackage{enumitem}

\captionsetup{
  font=small,
  labelfont={bf,color=azuldoc},
  labelsep=period,
  skip=6pt
}

% Gráficos vectoriales: TikZ, CircuiTikZ y PGFPlots
\usepackage{tikz}
\usetikzlibrary{
  babel,
  arrows.meta,
  positioning,
  shapes.geometric,
  shapes.multipart,
  fit,
  calc,
  automata,
  patterns,
  decorations.pathreplacing
}
\usepackage[american]{circuitikz}
\usepackage{pgfplots}
\pgfplotsset{compat=1.18}

% Pseudocódigos
\usepackage[ruled,vlined,linesnumbered,spanish]{algorithm2e}
\SetKwInput{KwDatos}{Entrada}
\SetKwInput{KwResultado}{Salida}
\SetAlCapSkip{0.6em}

% Código fuente (Python / C++ / protocolo)
\usepackage{listings}
\lstset{
  basicstyle=\ttfamily\footnotesize,
  backgroundcolor=\color{codigofondo},
  frame=single,
  rulecolor=\color{grisborde},
  breaklines=true,
  breakatwhitespace=true,
  tabsize=2,
  numbers=left,
  numberstyle=\tiny\color{gristexto},
  keywordstyle=\bfseries\color{azuldoc},
  commentstyle=\itshape\color{unalverde},
  stringstyle=\color{rojoborde},
  showstringspaces=false,
  captionpos=b,
  inputencoding=utf8,
  extendedchars=true,
  literate=
    {á}{{\'a}}1 {é}{{\'e}}1 {í}{{\'i}}1 {ó}{{\'o}}1 {ú}{{\'u}}1
    {Á}{{\'A}}1 {É}{{\'E}}1 {Í}{{\'I}}1 {Ó}{{\'O}}1 {Ú}{{\'U}}1
    {ñ}{{\~n}}1 {Ñ}{{\~N}}1 {ü}{{\"u}}1 {Ü}{{\"U}}1
    {°}{{\textdegree}}1 {±}{{$\pm$}}1 {×}{{$\times$}}1
    {≤}{{$\le$}}1 {≥}{{$\ge$}}1 {→}{{$\rightarrow$}}1 {—}{{---}}1
}

% Cajas semánticas (tcolorbox)
\usepackage[most]{tcolorbox}

\newtcolorbox{porcompletar}[1][Contenido pendiente por completar]{
  enhanced,
  breakable,
  colback=ambarfondo,
  colframe=ambarborde,
  coltitle=white,
  fonttitle=\bfseries\small,
  title={[POR COMPLETAR] #1},
  left=8pt,right=8pt,top=6pt,bottom=6pt,
  arc=3pt,
  boxrule=0.9pt
}

\newtcolorbox{decision}[1]{
  enhanced,
  breakable,
  colback=azulfondo,
  colframe=azuldoc,
  coltitle=white,
  fonttitle=\bfseries\small,
  title={Decisión de diseño --- #1},
  left=8pt,right=8pt,top=6pt,bottom=6pt,
  arc=3pt,
  boxrule=0.8pt
}

\newtcolorbox{hallazgo}[1][Hallazgo empírico]{
  enhanced,
  breakable,
  colback=verdefondo,
  colframe=unalverde,
  coltitle=white,
  fonttitle=\bfseries\small,
  title={#1},
  left=8pt,right=8pt,top=6pt,bottom=6pt,
  arc=3pt,
  boxrule=0.8pt
}

\newcommand{\placeholderfoto}[3][5.2cm]{%
  \begin{tcolorbox}[
    enhanced,
    colback=codigofondo,
    colframe=ambarborde,
    borderline={0.8pt}{2pt}{ambarborde,dashed},
    height=#1,
    valign=center,
    halign=center,
    arc=3pt
  ]
    \textbf{\color{ambarborde}[FOTOGRAFÍA POR INSERTAR]}\\[4pt]
    {\small\textbf{#2}}\\[3pt]
    {\footnotesize\color{gristexto}#3}
  \end{tcolorbox}
}

% Macros semánticas del proyecto
\newcommand{\cmd}[1]{\texttt{#1}}
\newcommand{\archivo}[1]{\texttt{#1}}
\newcommand{\estado}[3]{\ensuremath{(#1,#2,#3)}}
\newcommand{\badgeok}{\textcolor{unalverde}{\textbf{[Validado en pista / main]}}}
\newcommand{\badgepr}{\textcolor{azulclaro}{\textbf{[Implementado en rama / en integración]}}}
\newcommand{\badgepend}{\textcolor{ambarborde}{\textbf{[Planificado / pendiente]}}}

% Formato de títulos de sección
\titleformat{\section}
  {\Large\bfseries\color{azuldoc}}
  {\thesection.}{0.6em}{}
  [\vspace{-0.3em}\color{grisborde}\titlerule]

\titleformat{\subsection}
  {\large\bfseries\color{azuldoc}}
  {\thesubsection}{0.5em}{}

\titleformat{\subsubsection}
  {\normalsize\bfseries\color{gristexto}}
  {\thesubsubsection}{0.5em}{}

% Encabezado y pie de página
\pagestyle{fancy}
\fancyhf{}
\fancyhead[L]{\footnotesize\color{gristexto}Agente Solucionador de Laberintos 2D --- Inteligencia Artificial}
\fancyhead[R]{\footnotesize\color{gristexto}Universidad Nacional de Colombia}
\fancyfoot[C]{\footnotesize\color{gristexto}\thepage}
\renewcommand{\headrulewidth}{0.4pt}
\renewcommand{\footrulewidth}{0.2pt}

% Bibliografía (biblatex + biber)
\usepackage{csquotes}
\usepackage[backend=biber,style=ieee,sorting=none,urldate=comp]{biblatex}
\addbibresource{referencias.bib}

% Hipervínculos y referencias cruzadas (siempre al final del preámbulo)
\usepackage[
  colorlinks=true,
  linkcolor=azuldoc,
  citecolor=unalverde,
  urlcolor=azulclaro,
  plainpages=false,
  pdfpagelabels=true,
  pdftitle={Agente Autónomo Solucionador de Laberintos 2D en ESP32-S3},
  pdfauthor={Equipo Agente Laberinto - Universidad Nacional de Colombia}
]{hyperref}
\usepackage[spanish,nameinlink]{cleveref}
\crefname{algocf}{Algoritmo}{Algoritmos}
\Crefname{algocf}{Algoritmo}{Algoritmos}
\crefname{proposicion}{Proposición}{Proposiciones}
\Crefname{proposicion}{Proposición}{Proposiciones}
\crefname{definicion}{Definición}{Definiciones}
\Crefname{definicion}{Definición}{Definiciones}
